\section{From Error to Hazard}

The sections so far are strong on \emph{where errors originate} and \emph{why verification is hard}. What they don't yet supply is a way to connect an epistemic problem --- something wrong, missing, or uncertain in an output --- to an actual downstream loss. That connection needs its own vocabulary and its own analytical direction, run in addition to everything before it, not instead of it.

A vocabulary, stated as definitions rather than claims, since precision here prevents two specific bad habits. This report uses these terms as follows:

\begin{itemize}
\item \textbf{Error} --- a demonstrable departure from an applicable reference.
\item \textbf{Uncertainty} --- insufficient evidence to determine the answer either way.
\item \textbf{Ambiguity} --- more than one interpretation is reasonably available.
\item \textbf{Disagreement} --- competent evaluators reach different conclusions from the same evidence.
\item \textbf{Omission} --- required or relevant content was never considered or produced.
\item \textbf{Control violation} --- a required process or authorization wasn't followed, independent of whether the eventual result happened to be correct.
\item \textbf{Hazard} --- a system state capable of producing harm, whether or not harm has actually occurred yet.
\item \textbf{Incident} --- harm, or a defined near miss, has actually occurred.
\end{itemize}

This vocabulary matters most precisely in the weak-verification domains Section 7 already identified. Where no agreed reference answer exists, "error rate" is not a meaningful measurement --- what's actually being measured may be disagreement, incompleteness, instability, or noncompliance, each of which calls for a different response than "the model was wrong." Keeping these separate heads off two conclusions that are individually tempting and jointly dangerous: \textbf{"we found no factual errors, so the output is safe"} ignores omissions and control violations entirely --- an output can be factually error-free and still miss the one authority that mattered, or bypass a required approval step, and neither failure shows up in a check built only to catch false claims. And \textbf{"experts disagree, so verification is impossible"} confuses the absence of a single binary right-answer with the absence of any useful evidence at all --- disagreement among qualified reviewers is itself informative (it tells you the task sits toward the "expert assessor" or "unresolved ground truth" end of Section 7's verification-strength gradient), not a reason to stop checking.

\textbf{A hazard-first view, to complement rather than replace the five-function view.} Section 4's five functions begin with what a system \emph{does}. A hazard-first analysis begins instead with what \emph{must not happen}, and works backward from there. For a legal workflow, for instance, unacceptable losses might include: privileged material disclosed to the wrong party; a filing made against an incorrect deadline; controlling authority omitted from an analysis; unsupported advice communicated to a client as though it were settled; a matter barrier crossed; an irreversible action taken without proper authority. Starting from each of these, the analysis runs backward through five questions: What hazardous system state could produce this specific loss? What combination of model behavior, software behavior, data quality, human action, and organizational process could create that state? Which controls are meant to prevent it? What evidence actually shows those controls work, as opposed to being assumed to work? And what happens if they fail anyway?

This surfaces something the output-centered ladder in Section 9 structurally cannot: a completely correct model output can still participate in an unsafe outcome. It might be delivered to the wrong person, used well outside its intended scope, acted upon after it's gone stale, combined with other data that has since changed, or treated as authoritative by someone who lacked the standing to rely on it that way. None of that is a defect in the output itself --- the five-axis ladder would find nothing wrong with it --- and all of it can still produce real harm. This kind of backward-from-loss analysis has real precedent in established engineering practice: failure modes and effects analysis (FMEA) works upward from a component-level failure to ask how it affects the wider system, and systems-theoretic approaches such as STAMP go further, examining unsafe \emph{interactions} and inadequate \emph{control} even where no individual component has behaved in a way anyone would call, in isolation, a "failure."
